% ---------------------------------------------------------------------
\section{$\order{\as}$：实—虚相消与两圈关联函数}\label{sec:ee-oneloop}

本节在 $\order{\as}$ 用两种独立方法验证\cref{thm:ee-R-finite}：割图方法（虚修正与实辐射各含 $1/\varepsilon^2$ 极点，只有其和有限，\cref{prop:ee-oneloop}）与欧氏方法（两圈关联函数，从不出现红外极点，\cref{prop:ee-twoloop}）；二者在任意 $\varepsilon$ 下逐项相等（\cref{prop:ee-cutting-exact}）。所涉及的图见\cref{fig:ee-oneloop}。

\begin{proposition}[$\order{\as}$ 的实辐射与虚修正 \tagR]\label{prop:ee-oneloop}
（$R_{\mathrm{PT}}$：在重整化微扰论中对红外正规化的部分子态严格。）（单圈，即 $\order{\as}$。）用\cref{eq:ee-rho}的 $\rho$ 与公共归一化 $N_\varepsilon\equiv\frac{\as C_F}{2\pi}\big(\frac{4\pi\mu_0^2}{s}\big)^{\varepsilon}c_\Gamma$，两部分子割线上的虚修正与三部分子割线上的实辐射分别为
\begin{align}
\frac{\rho_V}{\rho_0}&=2\Real F^{(1)}(s+\ii0)=N_\varepsilon\cos(\pi\varepsilon)\,f(\varepsilon)=N_\varepsilon\Big[-\frac2{\varepsilon^2}-\frac3\varepsilon-8+\pi^2+\order{\varepsilon}\Big], \label{eq:ee-virt}\\
\frac{\rho_R}{\rho_0}&=\frac{\as C_F}{2\pi}\Big(\frac{4\pi\mu_0^2}{s}\Big)^{\varepsilon}\frac{2(1-2\varepsilon)(2-2\varepsilon+\varepsilon^2)\,\Gamma^2(1-\varepsilon)}{\varepsilon^2(1-3\varepsilon)(2-3\varepsilon)\,\Gamma(1-3\varepsilon)}\notag\\
&=N_\varepsilon\Big[\frac2{\varepsilon^2}+\frac3\varepsilon+\frac{19}2-\pi^2+\order{\varepsilon}\Big], \label{eq:ee-real}
\end{align}
% VF-EE-4, VF-EE-5
其中 $f(\varepsilon)\equiv-\frac{2-\varepsilon+2\varepsilon^2}{\varepsilon^2(1-2\varepsilon)}$；两式的闭式精确到 $\varepsilon$ 的一切阶。其和
\begin{equation}\label{eq:ee-R1}
\frac{\rho_V+\rho_R}{\rho_0}=N_\varepsilon\Big[\frac32+\order{\varepsilon}\Big]\quad\Longrightarrow\quad R=R_0\Big(1+\frac{3C_F}4\,\frac{\as}\pi\Big)+\order{\as^2},
\end{equation}
% VF-EE-6
SU(3) 中即 $R=R_0(1+\as/\pi)$。\emph{警告：}实、虚两部分各自的常数依赖于公共归一化：两种 $\Gamma$ 前因子相差 $\Gamma^2(1-\varepsilon)/\Gamma(1-3\varepsilon)=c_\Gamma\big(1-\frac{\pi^2}2\varepsilon^2+\order{\varepsilon^3}\big)$，恰好把实辐射的常数 $\frac{19}2$ 变成 $\frac{19}2-\pi^2$；只有同一归一化下的和才与归一化无关。（本命题只涉及部分子态。）
\uses{def:ee-R,prop:ee-R-ImPi,prop:sudakov-oneloop,def:prelim-dr-conventions,lem:appC-phase-space}
\end{proposition}
\begin{proof}
\emph{(1) 记账。}$\order{\as}$ 的 $\rho$ 有两类贡献（\cref{fig:ee-oneloop}）：$q\bar q$ 末态上 Born 振幅与单圈顶点的干涉 $\rho_V$，以及 $q\bar qg$ 末态的实辐射 $\rho_R$；无质量外线的自能是无标度积分，为零（\cref{def:prelim-dr-conventions}）。由\cref{prop:ee-R-ImPi}与\cref{eq:ee-Rratio}，$aR^{(1)}=R_0\lim_{\varepsilon\to0}(\rho_V+\rho_R)/\rho_0$；这一阶不需要耦合常数重整化。

\emph{(2) 虚修正。}对无质量夸克，手征守恒只允许 $\langle q\bar q|j^\mu|0\rangle=e_qF(q^2)\bar u\gamma^\mu v$，故 $\rho_V=2\Real F^{(1)}\rho_0$。由\cref{prop:sudakov-oneloop}（\tagR，$R_{\mathrm{PT}}$；推导见\cref{sec:appC-formfactor}），$F^{(1)}=-\frac{\as C_F}{4\pi}\big(\frac{4\pi\mu_0^2}{-q^2-\ii0}\big)^{\varepsilon}c_\Gamma\frac{2-\varepsilon+2\varepsilon^2}{\varepsilon^2(1-2\varepsilon)}$；$q^2=s+\ii0$ 处 $(-s-\ii0)^{-\varepsilon}=s^{-\varepsilon}\ee^{\ii\pi\varepsilon}$，故 $2\Real F^{(1)}=N_\varepsilon\cos(\pi\varepsilon)f(\varepsilon)$。由 $\cos\pi\varepsilon=1-\frac{\pi^2}2\varepsilon^2+\cdots$ 与 $f=-\frac2{\varepsilon^2}-\frac3\varepsilon-8-16\varepsilon+\cdots$ 得\cref{eq:ee-virt}：$+\pi^2$ 来自 $(-\frac2{\varepsilon^2})(-\frac{\pi^2}2\varepsilon^2)$，是类时延拓的效应；发散的相位不进入 $|F|^2$。

\emph{(3) 实辐射。}$\gamma^*(q)\to q(p_1)\bar q(p_2)g(k)$ 的振幅为
\[
\mathcal M^\nu=-g_0e_qt^a\,\bar u(p_1)\Big[\frac{\gamma^\alpha(\sla p_1+\sla k)\gamma^\nu}{s_{13}}-\frac{\gamma^\nu(\sla p_2+\sla k)\gamma^\alpha}{s_{23}}\Big]v(p_2)\,\varepsilon^*_\alpha(k),\qquad s_{ij}=2p_i\cdot p_j,\ p_3\equiv k .
\]
由 $\bar u(p_1)\sla k(\sla p_1+\sla k)=s_{13}\bar u(p_1)$ 与 $(\sla p_2+\sla k)\sla kv(p_2)=s_{23}v(p_2)$ 得 $k_\alpha\mathcal M^{\alpha\nu}=0$，故胶子极化求和取 $-g_{\alpha\beta}$ 是精确的。令 $y_{ij}=s_{ij}/s$（$y_{12}+y_{13}+y_{23}=1$），$D$ 维的迹给出
\begin{equation}\label{eq:ee-realme}
-g_{\mu\nu}\sum_{\rm pol,col}\mathcal M^\mu\mathcal M^{\nu*}=8(1-\varepsilon)g_0^2C_FN_ce_q^2\Big[(1-\varepsilon)\Big(\frac{y_{13}}{y_{23}}+\frac{y_{23}}{y_{13}}\Big)+\frac{2y_{12}}{y_{13}y_{23}}-2\varepsilon\Big] ,
\end{equation}
四维时即 $\frac{1}{\sigma_0}\frac{\dd\sigma}{\dd x_1\dd x_2}=\frac{C_F\as}{2\pi}\frac{x_1^2+x_2^2}{(1-x_1)(1-x_2)}$，$x_i=1-y_{jk}$。由\cref{lem:appC-phase-space}，三体相空间与单纯形上的 Dirichlet 积分为
\[
\dd\Phi_3=\frac{(4\pi)^{2\varepsilon-3}s^{1-2\varepsilon}}{2\Gamma(2-2\varepsilon)}\,\dd y_{13}\dd y_{23}\,(y_{12}y_{13}y_{23})^{-\varepsilon},\qquad
\int[\dd y]\,y_{12}^{a-1}y_{13}^{b-1}y_{23}^{c-1}=\frac{\Gamma(a)\Gamma(b)\Gamma(c)}{\Gamma(a+b+c)} .
\]
方括号三类项依次取 $(a,b,c)=(1-\varepsilon,2-\varepsilon,-\varepsilon)$（及对称者）、$(2-\varepsilon,-\varepsilon,-\varepsilon)$、$(1-\varepsilon,1-\varepsilon,1-\varepsilon)$；用 $\Gamma(2-\varepsilon)=(1-\varepsilon)\Gamma(1-\varepsilon)$、$\Gamma(-\varepsilon)=-\Gamma(1-\varepsilon)/\varepsilon$、$\Gamma(3-3\varepsilon)=(2-3\varepsilon)\Gamma(2-3\varepsilon)$ 及 $(1-\varepsilon)(2-3\varepsilon)-\varepsilon(1-\varepsilon)^2-\varepsilon^3=(1-2\varepsilon)(2-2\varepsilon+\varepsilon^2)$，
\[
I(\varepsilon)\equiv\int[\dd y]\,(y_{12}y_{13}y_{23})^{-\varepsilon}\big[\cdots\big]=\frac{2\Gamma^3(1-\varepsilon)}{\varepsilon^2\Gamma(3-3\varepsilon)}\,(1-2\varepsilon)(2-2\varepsilon+\varepsilon^2).
\]
除以 $\rho_0=N_ce_q^2(1-\varepsilon)s(4\pi/s)^\varepsilon\Gamma(1-\varepsilon)/[2\pi\Gamma(2-2\varepsilon)]$ 得 $\frac{\rho_R}{\rho_0}=8\pi g_0^2C_F\frac{(4\pi)^{\varepsilon-3}s^{-\varepsilon}}{\Gamma(1-\varepsilon)}I(\varepsilon)$，即\cref{eq:ee-real}的闭式。展开：$\frac{2(1-2\varepsilon)(2-2\varepsilon+\varepsilon^2)}{(1-3\varepsilon)(2-3\varepsilon)}=2+3\varepsilon+\frac{19}2\varepsilon^2+\order{\varepsilon^3}$，而由 $\ln\Gamma(1+x)=-\gamma_Ex+\frac{\pi^2}{12}x^2+\order{x^3}$，$\Gamma^2(1-\varepsilon)/\Gamma(1-3\varepsilon)=c_\Gamma(1-\frac{\pi^2}2\varepsilon^2+\cdots)$。

\emph{(4) 求和。}$1/\varepsilon^2$ 与 $1/\varepsilon$ 极点精确相消，有限部分 $-8+\pi^2+\frac{19}2-\pi^2=\frac32$，故 $aR^{(1)}/R_0=\frac{\as C_F}{2\pi}\cdot\frac32=\frac{3C_F}4\frac{\as}\pi$。警告由 (3) 的展开直接可见。
\end{proof}
这就是 Appelquist--Georgi\cite{AppelquistGeorgi1973}与 Zee\cite{Zee1973}的 $R=R_0(1+\as/\pi)$；教科书中的推导见\cite{EllisStirlingWebber1996}。上述每一步我们都用独立编写的计算机代数程序精确到 $\varepsilon$ 的一切阶核对过。

\begin{example}[KLN 的作用与局限]\label{ex:ee-quark-spectrum}
\cref{prop:ee-oneloop}的两项正是遍举截面 $\sigma(q\bar q)^{(1)}/\sigma_0$ 与 $\sigma(q\bar qg)/\sigma_0$，% VF-EE-11
各含 $\pm2/\varepsilon^2$ 的软—共线极点，只有对简并末态求和才有限：这是 Kinoshita--Lee--Nauenberg 定理\cite{Kinoshita1962,LeeNauenberg1964}（\cref{thm:ir-kln}，\tagA）的最简单实例。初态无色，故只需对\emph{末态}求和；有色初态时非阿贝尔 Bloch--Nordsieck 相消的失效\cite{DoriaFrenkelTaylor1980,DiLietoEtAl1981}与此无关（一般过程中对初态或末态之一的简并态求和即可，但通常须计入前向散射以及不连通图与连通图的干涉\cite{Frye:2018xjj}）。再看一个单举但不红外安全的量：夸克能谱 $\dd\sigma/\dd x$，$x=2E_q/\sqrt s=1-y_{23}$。由\cref{eq:ee-realme}与 $\dd\Phi_3$，$\frac1{\rho_0}\frac{\dd\rho_R}{\dd y_{13}\dd y_{23}}=\frac{\as C_F}{2\pi}\big(\frac{4\pi\mu_0^2}{s}\big)^\varepsilon\frac{(y_{12}y_{13}y_{23})^{-\varepsilon}}{\Gamma(1-\varepsilon)}[\cdots]$。固定 $x<1$ 时 $y_{13}\in[0,x]$ 包含夸克与胶子共线的区域，那里方括号趋于 $y_{13}^{-1}[(1-\varepsilon)(1-x)+\frac{2x}{1-x}]$，而 $\int_0\dd y_{13}\,y_{13}^{-1-\varepsilon}=-\frac{x^{-\varepsilon}}\varepsilon+\cdots$，于是
\begin{equation}\label{eq:ee-xspec}
\frac1{\sigma_0}\frac{\dd\sigma}{\dd x}\Big|_{\text{极点}}=-\frac1\varepsilon\,\frac{\as}{2\pi}\,C_F\,\frac{(1-x)^2+2x}{1-x}=-\frac1\varepsilon\,\frac{\as}{2\pi}\,P_{qq}(x)\qquad(x<1),
\end{equation}
$P_{qq}=C_F[\frac{1+x^2}{1-x}]_+$。连同 $x=1$ 处的虚修正与软端点，矩 $\int_0^1x^n\frac{\dd\sigma}{\dd x}$ 的 $1/\varepsilon^2$ 都相消，$1/\varepsilon$ 的系数（单位 $\frac{\as}{2\pi}$）为 $-\int_0^1x^nP_{qq}=0,\frac43C_F,\frac{25}{12}C_F,\frac{157}{60}C_F$（$n=0,\dots,3$；例如 $n=1$ 时 $-C_F\int_0^1(x-1)\frac{1+x^2}{1-x}=\frac43C_F$）。$n=0$ 是总截面，由夸克数守恒而有限；$n\ge1$ 的剩余共线极点由后文定义的 $\overline{\rm MS}$ 碎裂函数（\cref{def:pdf-ff}）吸收，这就是后文单举湮灭的因子化（\cref{thm:coll-sia}，\tagA）的出发点。KLN 只保护对全部简并态求和的量。
\uses{prop:ee-oneloop,thm:ir-kln}
\end{example}

%%FIG:fig_ee_oneloop%%

\begin{proposition}[欧氏两圈关联函数 \tagR]\label{prop:ee-twoloop}
（$R_{\mathrm{PT}}$：在重整化微扰论中对红外正规化的部分子态严格。）（两圈，即 $\order{g_0^2}$。）对一味夸克，流关联函数的两圈裸贡献精确到 $\varepsilon$ 的一切阶为
\begin{equation}\label{eq:ee-Pi2bare}
\Pi^{(1)}_{\rm bare}=N_ce_q^2C_F\,\frac{g_0^2(Q^2)^{-2\varepsilon}}{(4\pi)^D}\Big[r_1\,G(1,1)^2+r_2\,G(1,1)\,G(1+\varepsilon,1)\Big],\qquad Q^2=-q^2-\ii0,
\end{equation}
\begin{gather*}
r_1=-\frac{4(1-\varepsilon)(2-\varepsilon+2\varepsilon^2)}{\varepsilon(3-2\varepsilon)},\qquad
r_2=\frac{16(1-\varepsilon)(1-2\varepsilon)(2-2\varepsilon+\varepsilon^2)}{\varepsilon(3-2\varepsilon)(2-3\varepsilon)},\\
G(1+\varepsilon,1)=\frac{\Gamma(2\varepsilon)\Gamma(1-2\varepsilon)\Gamma(1-\varepsilon)}{\Gamma(1+\varepsilon)\Gamma(2-3\varepsilon)} ;
\end{gather*}
% VF-EE-7
它满足 $q_\mu q_\nu\Pi^{\mu\nu}=0$（精确到 $\varepsilon$ 的一切阶）。$\overline{\rm MS}$ 重整化后，对全部夸克味求和，
\begin{equation}\label{eq:ee-PiMSbar}
\Pi_{\overline{\rm MS}}(q^2)=\frac{R_0}{12\pi^2}\Big[\frac53-L+a\,\frac{3C_F}4\Big(\frac{55}{12}-4\zeta_3-L\Big)\Big]+\order{a^2},\qquad
D(Q^2)=R_0\Big(1+\frac{3C_F}4\,a\Big)+\order{a^2} .
\end{equation}
常数 $\frac53$ 与 $\frac{55}{12}-4\zeta_3$ 依赖于方案（接触项），$L$ 的系数不依赖。（本命题只涉及部分子态的 Green 函数。）
\uses{def:prelim-msbar,def:prelim-feynman-rules,lem:appC-bubble,lem:appC-ibp}
\end{proposition}
\begin{proofsketch}
(1)~\emph{图与因子。}$\order{g_0^2}$ 有三个图（\cref{fig:ee-oneloop}(d)--(f)）：胶子连接两条夸克线的 V，自能插入 S$_1$、S$_2$。由\cref{def:prelim-feynman-rules}（Feynman 规范），每个图带因子 $\ii\cdot(-1)\cdot(\ii g_0)^2\cdot(-\ii)\cdot\ii^4\cdot C_FN_c=g_0^2C_FN_c$（依次来自 $\Pi^{\mu\nu}$ 的定义、费米子圈、两个顶点、胶子与四个夸克传播子）。取 $r=k+l$，Wick 转动后分母属于两圈传播子拓扑
\begin{equation}\label{eq:ee-topo}
\begin{gathered}
D_1=k^2,\quad D_2=(k-q)^2,\quad D_3=r^2,\quad D_4=(r-q)^2,\quad D_5=(k-r)^2;\\
\mathrm V:(1,1,1,1,1),\qquad \mathrm S_1:(2,1,1,0,1),\qquad \mathrm S_2:(1,2,0,1,1)
\end{gathered}
\end{equation}
（括号内为各分母的幂次）。
(2)~\emph{分子。}用 $-g_{\mu\nu}$（得 $\Pi$）与 $q_\mu q_\nu$（检验 Ward 恒等式）投影，算出 $D$ 维的迹；五个标量积都是 $D_i$ 与 $Q^2$ 的线性组合，每个图化为 $\hat I(a)\equiv\int\!\!\int\frac{\dd^Dk\,\dd^Dr}{\pi^D}\prod_iD_i^{-a_i}$（欧氏，乘以 $(Q^2)^{D-\sum a_i}$）之和。
(3)~\emph{积分。}$a_5\le0$ 时积分分解为两个单圈积分之积（\cref{lem:appC-bubble}及其张量公式）；$a_1,\dots,a_4$ 中某个 $\le0$ 时先做的单圈子积分给出非整数幂的传播子，再由 $G(n,m+\varepsilon)$ 完成；全为正时用 Chetyrkin--Tkachov 的三角规则\cite{ChetyrkinTkachov1981,Tkachov1981}（\cref{lem:appC-ibp}）降低 $a_3$、$a_4$ 或 $a_5$；例如 $\hat I(1,1,1,1,1)=6\zeta_3+\order{\varepsilon}$。
(4)~\emph{结果。}全部积分化为两种超越结构 $G(1,1)^2$ 与 $G(1,1)G(1+\varepsilon,1)$，即\cref{eq:ee-Pi2bare}；$q_\mu q_\nu\Pi^{\mu\nu}=0$ 精确成立。$r_1$ 含形状因子的多项式 $2-\varepsilon+2\varepsilon^2$，$r_2$ 含实辐射的多项式 $2-2\varepsilon+\varepsilon^2$，分别对应两粒子割与三粒子割（下文的切割恒等式把这一对应精确化）。
(5)~\emph{逐图的极点。}乘以 $\mu^{2\varepsilon}S_\varepsilon^{-1}$，以 $aC_F\frac{N_ce_q^2}{12\pi^2}$ 为单位（Feynman 规范）：
\[
\begin{aligned}
\mathrm V&=\tfrac{1}{4\varepsilon^2}+\tfrac{1}{\varepsilon}\big(\tfrac{17}{12}-\tfrac L2\big)+\tfrac{L^2}2-\tfrac{17}6L-3\zeta_3-\tfrac{\pi^2}{24}+\tfrac{259}{36},\\
\mathrm S_1+\mathrm S_2&=-\tfrac{1}{4\varepsilon^2}+\tfrac1\varepsilon\big(\tfrac L2-\tfrac{25}{24}\big)-\tfrac{L^2}2+\tfrac{25}{12}L+\tfrac{\pi^2}{24}-\tfrac{541}{144},
\end{aligned}
\]
和为 $\frac{3}{8\varepsilon}+\frac{55}{16}-3\zeta_3-\frac34L+\order{\varepsilon}$：$1/\varepsilon^2$ 与非局域的 $L/\varepsilon$ 相消（子发散按 Ward 恒等式 $Z_1=Z_2$ 相消），剩下的局域单极点由 $Z_\Pi$ 吸收。
(6)~\emph{$\overline{\rm MS}$。}单圈：$12\pi^2\mu^{2\varepsilon}S_\varepsilon^{-1}\Pi_0/(N_ce_q^2)=\frac1\varepsilon+\frac53-L+\order{\varepsilon}$；两圈：去掉 $\frac3{8\varepsilon}$ 后剩 $\frac34(\frac{55}{12}-4\zeta_3-L)$。对各味求和即得\cref{eq:ee-PiMSbar}，$D=-12\pi^2\partial_L\Pi_{\overline{\rm MS}}$。完整论证见\cref{sec:appG-ee}；我们用独立编写的计算机代数程序（自编的 $D$ 维迹、张量公式与分部积分递推）逐图重复了这一计算，结果精确一致。
\end{proofsketch}
两圈的局域极点也给出 $\gamma_\Pi=\frac{R_0}{12\pi^2}(1+\frac{3C_F}4a)+\order{a^2}$（\cref{prop:ee-adler-rg}(a)）。

\begin{example}[欧氏两圈路线]\label{ex:ee-twoloop-route}
按\cref{thm:ee-R-finite}(iii)，在\cref{eq:ee-PiMSbar}中代入 $L\to\ell-\ii\pi$：$\Imag(\frac53-L)=\pi$，$\Imag[\frac{3C_F}4(\frac{55}{12}-4\zeta_3-L)]=\frac{3C_F}4\pi$，故 $\Imag\Pi(s+\ii0)=\frac{R_0}{12\pi}(1+\frac{3C_F}4a)$，$R=R_0(1+\frac{3C_F}4\frac{\as}\pi)$，与\cref{eq:ee-R1}一致。这是\cref{thm:ee-R-finite}最直接的演示：裸的两圈 $\Pi$ 只有一个局域的（紫外的）单极点，依赖方案的常数是接触项，在 $R$ 中消失，$R$ 只由 $L$ 的系数（即 $D$）决定；割图方法中 $\pm2/\varepsilon^2$ 的红外极点在这里根本不出现。
\uses{prop:ee-twoloop,thm:ee-R-finite}
\end{example}

\begin{proposition}[$O(g^2)$ 处精确到 $\varepsilon$ 的切割恒等式 \tagR]\label{prop:ee-cutting-exact}
（$R_{\mathrm{PT}}$：在重整化微扰论中对红外正规化的部分子态严格。）（$\order{g_0^2}$ 阶。）对一切 $\varepsilon$（各 $\Gamma$ 函数的极点除外），
\begin{equation}\label{eq:ee-cutexact}
2(D-1)\,s\,\Imag\Pi^{(1)}_{\rm bare}(s+\ii0)=\rho_V+\rho_R ,
\end{equation}
% VF-EE-8
而且逐项成立：\cref{eq:ee-Pi2bare}中 $r_1G(1,1)^2$ 一项等于 $\rho_V$（两粒子割），$r_2G(1,1)G(1+\varepsilon,1)$ 一项等于 $\rho_R$（三粒子割），$\rho_{V,R}$ 取\cref{prop:ee-oneloop}的闭式（以 $g_0^2=4\pi\as\mu_0^{2\varepsilon}$ 表示）。（本命题只涉及部分子态。）
\uses{prop:ee-oneloop,prop:ee-twoloop}
\end{proposition}
\begin{proof}
记 $P\equiv N_ce_q^2C_Fg_0^2\,s\,(4\pi/s)^{2\varepsilon}/(4\pi)^4$，$G_{11}\equiv G(1,1)$。由 $(-s-\ii0)^{-2\varepsilon}=s^{-2\varepsilon}\ee^{2\pi\ii\varepsilon}$，
\[
2(D-1)\,s\,\Imag\Pi^{(1)}_{\rm bare}(s+\ii0)=2(3-2\varepsilon)P\sin(2\pi\varepsilon)\big[r_1G_{11}^2+r_2G_{11}G_{1+\varepsilon,1}\big],\qquad G_{11}=\frac{\Gamma(\varepsilon)\Gamma^2(1-\varepsilon)}{\Gamma(2-2\varepsilon)} .
\]反复用 $\Gamma(z)\Gamma(1-z)=\pi/\sin\pi z$（$z=\varepsilon,2\varepsilon$）与 $\sin2\pi\varepsilon=2\sin\pi\varepsilon\cos\pi\varepsilon$。另一方面，$\frac{\as}{2\pi}(\frac{4\pi\mu_0^2}{s})^\varepsilon=\frac{g_0^2}{8\pi^2}(\frac{4\pi}s)^\varepsilon$，$\rho_0=N_ce_q^2(1-\varepsilon)s(4\pi/s)^\varepsilon\Gamma(1-\varepsilon)/[2\pi\Gamma(2-2\varepsilon)]$。

\emph{第一项。}$2(3-2\varepsilon)r_1=-8(1-\varepsilon)(2-\varepsilon+2\varepsilon^2)/\varepsilon$，$G_{11}^2\sin2\pi\varepsilon=2\pi\cos(\pi\varepsilon)\Gamma(\varepsilon)\Gamma^3(1-\varepsilon)/\Gamma^2(2-2\varepsilon)$，故第一项为 $-16\pi P\cos(\pi\varepsilon)(1-\varepsilon)(2-\varepsilon+2\varepsilon^2)\Gamma(\varepsilon)\Gamma^3(1-\varepsilon)/[\varepsilon\Gamma^2(2-2\varepsilon)]$。而 $\rho_V=2\Real F^{(1)}\rho_0=16\pi P\,c_\Gamma\cos(\pi\varepsilon)f(\varepsilon)(1-\varepsilon)\Gamma(1-\varepsilon)/\Gamma(2-2\varepsilon)$；由 $c_\Gamma/[\varepsilon(1-2\varepsilon)]=\Gamma(\varepsilon)\Gamma^2(1-\varepsilon)/\Gamma(2-2\varepsilon)$，$c_\Gamma f=-(2-\varepsilon+2\varepsilon^2)\Gamma(\varepsilon)\Gamma^2(1-\varepsilon)/[\varepsilon\Gamma(2-2\varepsilon)]$，故 $\rho_V$ 恰等于第一项。

\emph{第二项。}$2(3-2\varepsilon)r_2=32(1-\varepsilon)(1-2\varepsilon)(2-2\varepsilon+\varepsilon^2)/[\varepsilon(2-3\varepsilon)]$，由 $\Gamma(2\varepsilon)\Gamma(1-2\varepsilon)\sin2\pi\varepsilon=\pi$ 与 $\Gamma(\varepsilon)/\Gamma(1+\varepsilon)=1/\varepsilon$，$G_{11}G_{1+\varepsilon,1}\sin2\pi\varepsilon=\pi\Gamma^3(1-\varepsilon)/[\varepsilon\Gamma(2-2\varepsilon)\Gamma(2-3\varepsilon)]$，第二项为 $32\pi P(1-\varepsilon)(1-2\varepsilon)(2-2\varepsilon+\varepsilon^2)\Gamma^3(1-\varepsilon)/[\varepsilon^2(2-3\varepsilon)\Gamma(2-2\varepsilon)\Gamma(2-3\varepsilon)]$；用 $(1-3\varepsilon)\Gamma(1-3\varepsilon)=\Gamma(2-3\varepsilon)$，$\rho_R=\frac{\rho_R}{\rho_0}\rho_0$ 恰为同一表达式。
\end{proof}
这是 $\order{g^2}$ 的幺正性在 $D$ 维中的直接验证：欧氏两圈积分与闵氏相空间积分在 $\varepsilon\to0$ \emph{之前}就一致，两种超越结构与两类割线一一对应（我们还在若干 $\varepsilon$ 值上以数十位精度作了数值核对）。这一逐项恒等式是本书给出的推导，依据是维数正规化中标准的幺正性\cite{Cutkosky1960,Veltman1963,tHooftVeltman1974}，不主张其新颖性。

\begin{corollary}[与部分子末态求和的等同；KLN \tagA]\label{cor:ee-cuts}
在 $n_f$ 味无质量 QCD 中，\cref{thm:ee-R-finite}的 $R^{(n)}$ 等于 $\gamma^*\to$ 部分子的全部割图之和，即对一切部分子末态（物理极化）求 $D$ 维截面之和，包括 $\overline{\rm MS}$ 抵消项的插入，再取 $\varepsilon\to0$。每一类割线单独含有直到 $\varepsilon^{-2n}$ 的极点，只有全部求和才有限；对 $e^+e^-$ 只需对末态的简并态求和（\cref{fig:ee-cutsum-disc}）。（本推论只涉及部分子态。）
\uses{thm:ee-R-finite,thm:prelim-cutting-rules,rem:prelim-cutting-dr,prop:prelim-unitarity,prop:ee-cutting-exact}
\end{corollary}
\begin{proofsketch}
(1)~对每个 $D$ 维无质量图，把两个流顶点分开的全部割线之和等于该图的不连续性 $2\Imag\Pi_{\rm graph}(s+\ii0)$：最大时间方程（\cref{thm:prelim-cutting-rules}，\tagR，$R_{\mathrm{PT}}$），其无质量维数正规化形式是\cref{rem:prelim-cutting-dr}（\tagA）。(2)~鬼场割线与非物理极化的割线相消（\cref{prop:prelim-unitarity}，\tagR，$R_{\mathrm{PT}}$），割线之和等于物理末态之和。(3)~$Z_\Pi$ 为实且与 $q$ 无关，没有不连续性；耦合常数抵消项对应割图中的抵消项顶点。(4)~由\cref{thm:ee-R-finite}，每一阶的和有限且等于 $R^{(n)}$；每发射一个部分子至多带来 $\varepsilon^{-2}$。(5)~$\order{g^2}$ 由\cref{prop:ee-cutting-exact}精确验证。割图规则见\cite{Cutkosky1960,Veltman1963,tHooftVeltman1974}，鬼场与幺正性见\cite{tHooftVeltman1972b,KugoOjima1979}；完整论证见\cite{Sterman1978a,Sterman1978b}。
\end{proofsketch}

\paragraph{地位与缺口}最弱的一步是无质量、以维数正规化为红外正规子的 $D$ 维割图规则（含鬼场）：最大时间方程与 BRST 幺正性各自是标准的，但没有一篇文献把这一情形写到完全严格（\cref{rem:prelim-cutting-dr}，\tagA）。这一缺口不是某个幂次估计的缺口，而是一个恒等式（亚纯函数之间的等式）的完整证明，所以误差类型的分级在这里不适用。它不影响\cref{thm:ee-R-finite}，那里 $R$ 由 $12\pi\Imag\Pi$ 定义（由整个理论的幺正性，它就是物理截面）。若把割图之和取作 $R^{(n)}$ 的定义，其有限性目前只有全阶论证；$\order{g^2}$ 的精确验证（\cref{prop:ee-cutting-exact}）与 Sterman 在闵氏空间中的全阶论证\cite{Sterman1978a,Sterman1978b}是证据。提升为 \tagR 需要维数正规化中无质量、含鬼场的割图规则的完整证明（例如由 BRST 不变的最大时间方程出发）。
